\documentclass[11pt,a4paper]{article}

\usepackage[margin=2.5cm]{geometry}
\usepackage{amsmath}
\usepackage{graphicx}
\usepackage{hyperref}
\usepackage{enumitem}
\usepackage{booktabs}
\usepackage{xcolor}
\usepackage{fancyhdr}
\usepackage{titlesec}

\hypersetup{
    colorlinks=true,
    linkcolor=blue!60!black,
    citecolor=blue!60!black,
    urlcolor=blue!60!black
}

\pagestyle{fancy}
\fancyhf{}
\fancyhead[L]{\small\textit{Architecture Overview}}
\fancyhead[R]{\small\textit{v0.2 -- March 2026}}
\fancyfoot[C]{\thepage}

\title{%
    \vspace{-1cm}
    \textbf{Breach}\\[0.3cm]
    \Large A Tactical Spaceship Boarding Game\\[0.2cm]
    \large Document 00: Architecture Overview \& Design Decisions
}
\author{Erik \and Idir}
\date{March 2026}

\begin{document}
\maketitle

\begin{abstract}
This document summarizes the core architectural decisions for
\textbf{Breach}, a tactical squad-based spaceship boarding game.  It
serves as the foundation reference for all subsequent system-specific
design documents and as a discussion basis for the development team.
\end{abstract}

\tableofcontents
\newpage

%----------------------------------------------------------------------
\section{Game Overview}
%----------------------------------------------------------------------

A top-down, turn-based tactical game in which squads board spaceships to
complete objectives (assassination, theft, hostage rescue, etc.).  The
core design philosophy is \textbf{emergent complexity from simple rules}:
a small number of interconnected systems---destructible environments,
pressure, fire, liquids, creatures---combine to produce a vast space of
possible situations.

The game is set in a satirical 2087 world ruled by central banking
dynasties.  Tone is darkly humorous and absurdist.

%----------------------------------------------------------------------
\section{Engine \& Tooling}
\label{sec:engine}
%----------------------------------------------------------------------

The core architecture is a \textbf{headless simulation layer} fully
separated from the presentation layer (Section~\ref{sec:twolayer}).
This means the engine choice affects \emph{only} the presentation
layer---rendering, input, audio, UI---while the simulation remains
engine-agnostic.

\subsection{Engine Candidates}

\begin{description}[style=nextline]
    \item[Unreal Engine]
        Idir's primary expertise.  Powerful rendering, mature tooling,
        C++ native.  However, Unreal is fundamentally 3D-first---heavy
        for a 2D tile-based game with significant overhead that goes
        unused.  Makes more sense if the project later moves to 2.5D
        or 3D.

    \item[Unity]
        Excellent 2D support with a built-in 2D lighting system that
        supports normal maps natively.  Large asset store and community
        for 2D game development.  Lighter weight than Unreal for 2D
        projects.  The simulation is C++ regardless; Unity calls it via
        native plugin.  Licensing terms should be reviewed.

    \item[Godot]
        Free, open source, lightweight, with excellent 2D support.
        Built-in \texttt{CanvasItem} shaders with normal map support.
        C++ integration via GDExtension.  Smaller ecosystem than
        Unity/Unreal but very well suited to 2D games with dynamic
        lighting.

    \item[Custom Engine]
        Total control, no licensing issues, educational value.  But an
        enormous time investment on solved problems (rendering, input,
        audio, UI, asset loading).  Not recommended unless engine
        development is a goal in itself.
\end{description}

\subsection{Recommendation}

For a 2D tile-based game with dynamic lighting and normal-mapped
sprites, \textbf{Unity or Godot} are the strongest candidates.  Both
have native 2D lighting pipelines that match our design requirements
(see Document~13: Graphics, Lighting \& Art Pipeline).  Unreal remains
viable---especially given Idir's expertise---but the 2D lighting
workflow may be smoother in Unity or Godot.

\medskip
\noindent\textbf{Decision status:} Open---to be discussed and confirmed
with Idir before prototyping begins.

%----------------------------------------------------------------------
\section{Two-Layer Architecture}
\label{sec:twolayer}
%----------------------------------------------------------------------

\begin{description}[style=nextline]
    \item[Simulation Layer]
        A self-contained simulation that owns all game rules and state.
        This layer has \textbf{no dependency on any game engine} and can
        run headless for AI training and automated testing.  It exposes
        a clean interface: \texttt{get\_state()},
        \texttt{apply\_action()}, \texttt{step()}.

    \item[Presentation Layer]
        The game engine reads simulation state each frame and renders
        it.  Player input is captured by the engine, translated into
        game actions, and forwarded to the simulation.

    \item[Neural Network Training]
        The headless simulation generates training data and runs
        self-play at scale without spinning up a rendering engine.
\end{description}

This separation allows parallel development: simulation logic can
progress independently of visual implementation.

%----------------------------------------------------------------------
\section{Spatial Representation}
\label{sec:spatial}
%----------------------------------------------------------------------

\subsection{Grid-Based World}

The game world is represented as a \textbf{uniform 2D grid}.  Every
structural element (walls, floors, hull, doors, windows) occupies one or
more grid cells.  This was chosen over continuous geometry for the
following reasons:

\begin{itemize}
    \item All propagation systems (pressure, fire, liquid) become
          cellular automata---simple, fast, and predictable.
    \item Game state maps directly to a 2D tensor for neural network
          input (stacked feature planes, cf.\ AlphaStar).
    \item Pathfinding, line-of-sight, and cover calculations are
          straightforward on a grid.
    \item Serialization, replay, and debugging are trivial.
\end{itemize}

\subsection{Fine Grid with Coarse Logical Structure}

A single player entity occupies a \textbf{3$\times$3 block} of tiles.
The grid resolution is therefore finer than the logical ``square'' a
character stands in.  This gives us:

\begin{itemize}
    \item Smoother movement and sub-positioning (center, edges, corners).
    \item Granular destruction---a wall segment can have a small hole,
          a large hole, or be fully destroyed.
    \item Natural partial cover and line-of-sight without a separate
          high-resolution system.
    \item Better visual fidelity for gore, splatter patterns, and debris.
\end{itemize}

Room layout, doors, and corridors are still designed on the
\textbf{coarse grid} (each coarse cell = 3$\times$3 fine tiles).
The fine grid provides resolution \emph{within} that structure.

\medskip
\noindent\textbf{Implementation note:}  Grid resolution should be a
configurable parameter.  Initial development can use 1:1 (one tile per
logical cell) and scale to 3:1 once core systems are stable.

%----------------------------------------------------------------------
\section{State Representation}
\label{sec:state}
%----------------------------------------------------------------------

\subsection{Tile Objects + Cached Arrays}

The map is a 2D grid of \textbf{Tile objects}---each tile stores its
material, hit points, fire state, liquid contents, and a list of
entities.  This object-based representation is easy to extend (adding a
property means adding a field, not changing every tensor consumer) and
naturally handles mixed state (type information alongside continuous
values like fire intensity).

For physics and propagation, the simulation maintains \textbf{cached
arrays} derived from tile state: light-blocking, gas-blocking,
walkability, flammability.  These caches are rebuilt on tile change
(\texttt{on\_tile\_changed(x, y)}) and allow all propagation systems to
operate as fast bulk array operations.

This hybrid gives us the best of both worlds: rich per-tile data for
game logic, and fast array operations for physics.

\subsection{Data-Driven Material System}

Tile behavior is determined by \textbf{material type}.  All material
properties are defined in a single lookup table---adding a new material
requires one table entry and zero code changes:

\begin{table}[h]
\centering
\begin{tabular}{@{}lccccc@{}}
\toprule
\textbf{Material} & \textbf{Max HP} & \textbf{Flammable} & \textbf{Blocks Light} & \textbf{Blocks Gas} & \textbf{Blast Resist.} \\
\midrule
None   & 0   & No  & No  & No  & 0.0 \\
Steel  & 200 & No  & Yes & Yes & 0.8 \\
Wood   & 60  & Yes & Yes & Yes & 0.2 \\
Glass  & 15  & No  & No  & Yes & 0.0 \\
Hull   & 300 & No  & Yes & Yes & 0.9 \\
\bottomrule
\end{tabular}
\caption{Material properties.  Each tile's behavior derives entirely
from its material.  The level editor paints materials; all properties
auto-derive.}
\label{tab:materials}
\end{table}

\subsection{State Layers}

The following fields are tracked per tile, either as Tile object
properties or as cached arrays for bulk computation:

\begin{table}[h]
\centering
\begin{tabular}{@{}lll@{}}
\toprule
\textbf{Field}       & \textbf{Storage}   & \textbf{Description} \\
\midrule
Material             & Tile object         & Determines all structural properties \\
Hit points           & Tile object         & Per-tile structural HP \\
Pressure             & Cached array        & Atmospheric pressure (gradual decompression) \\
Fire intensity       & Tile object         & Heat / flame level \\
Liquid type          & Tile object         & None, blood, water, fuel \ldots \\
Liquid depth         & Tile object         & Volume of liquid in tile \\
Smoke density        & Cached array        & Obscures vision, spreads through gaps \\
Oxygen               & Cached array        & Supports suffocation, consumed by fire \\
Temperature          & Tile object         & Thermal state \\
Entity list          & Tile object         & Creatures/players occupying tile \\
\bottomrule
\end{tabular}
\caption{Per-tile state fields.  ``Cached array'' fields are also
stored per-tile but maintained as contiguous arrays for physics
computation.  Additional layers may be added per system document.}
\label{tab:fields}
\end{table}

\subsection{Why This Approach?}

\begin{itemize}
    \item \textbf{Game logic as array operations.}  Explosions, pressure
          propagation, and fire spread operate on cached arrays---fast
          bulk computation via NumPy (prototype) or Eigen/raw arrays
          (C++).
    \item \textbf{Emergent interactions are automatic.}  Each system
          reads shared arrays and writes to its own.  Fire checks the
          liquid array for fuel; creatures check it for blood.  No
          explicit cross-system wiring.
    \item \textbf{ML-ready.}  Stack the cached arrays into channels and
          feed directly to a CNN---no conversion step.
    \item \textbf{Extensible.}  Adding a new tile property is a field
          addition, not a tensor schema change.
    \item \textbf{Serializable.}  Dump tile objects $\rightarrow$
          complete snapshot.  Save/load, replay, and undo come nearly
          for free.
\end{itemize}

\subsection{Entity List}

Mobile entities (squad members, creatures, items) are stored in a
\textbf{separate entity list}, since they carry properties that do not
map well to spatial grids: inventory, action points, status effects,
allegiance, AI state.  Their position is an $(x, y)$ index into the
grid.

%----------------------------------------------------------------------
\section{Simulation Tick \& Double Buffering}
\label{sec:tick}
%----------------------------------------------------------------------

\subsection{Turn Resolution}

Each player turn proceeds as follows:

\begin{enumerate}
    \item Player (or AI) selects and submits actions.
    \item Actions are applied as direct state modifications (damage,
          movement, item use, etc.).
    \item The simulation runs $N$ propagation ticks to let the
          environment settle (pressure equalization, fire spread, liquid
          flow, creature reactions).
    \item The new steady-state is presented; the next player's turn
          begins.
\end{enumerate}

\subsection{Double-Buffered Propagation}

All propagation systems resolve \textbf{simultaneously} within each
tick.  This is implemented via double buffering:

\begin{enumerate}
    \item Each cached array has a \emph{current} buffer (read-only
          during the tick) and a \emph{next} buffer (write target).
    \item Every system reads from \emph{current} and writes to
          \emph{next}.
    \item At the end of the tick, \emph{next} becomes \emph{current}
          (pointer swap).
\end{enumerate}

This avoids order-of-operations bugs and produces physically more
accurate behavior.

%----------------------------------------------------------------------
\section{Design Document Roadmap}
\label{sec:roadmap}
%----------------------------------------------------------------------

Each system is specified in its own design document.  Documents progress
through: \textbf{brainstorm} (Markdown notes) $\rightarrow$
\textbf{design document} (\LaTeX) $\rightarrow$ \textbf{implementation}.

\subsection{Environmental Simulation}

Matrix-based physics and propagation systems---the cellular automata
that drive emergent gameplay.

\begin{table}[h]
\centering
\begin{tabular}{@{}llp{5.5cm}l@{}}
\toprule
\textbf{Doc} & \textbf{System} & \textbf{Primary State (read/write)} & \textbf{Status} \\
\midrule
01 & Structure, Materials \& Destruction & Material (rw), HP (rw), Pressure (w) & Designed \\
02 & Pressure \& Decompression & Pressure (rw), Oxygen (rw), Gas-block (r) & Designed \\
03 & Fire Propagation & Fire (rw), Oxygen (r), Liquid (r), HP (w) & Brainstorm \\
04 & Liquids (blood, water, fuel) & Liquid type/depth (rw), Pressure (r) & Brainstorm \\
05 & Smoke Propagation & Smoke (rw), Pressure (r), Gas-block (r) & Designed \\
06 & Explosions & Pressure (w), HP (w), Fire (w), Smoke (w) & Designed \\
07 & Line of Sight \& Cover & Light-block (r), Fire (r), Smoke (r) & Brainstorm \\
\bottomrule
\end{tabular}
\caption{Environmental simulation documents.  ``Designed'' indicates
that core equations and algorithms have been worked out in brainstorm
notes; the \LaTeX{} design document is pending.}
\label{tab:env-systems}
\end{table}

\subsection{Game Mechanics}

Higher-level rules and systems that run on top of the environmental
simulation.

\begin{table}[h]
\centering
\begin{tabular}{@{}llp{5.5cm}l@{}}
\toprule
\textbf{Doc} & \textbf{System} & \textbf{Scope} & \textbf{Status} \\
\midrule
08 & Turn System \& Simultaneous Turns & Turn structure, simultaneous resolution, action points & Brainstorm \\
09 & Combat Mechanics & Damage, weapons, hit calculation & Brainstorm \\
10 & Economy \& Currency & CBD / Gold / Bitcoin, inflation, tracking & Brainstorm \\
11 & Inventory \& Equipment & Loadout, weight, item interactions & Brainstorm \\
12 & Player Interface \& Interaction & Controls, HUD, action selection, feedback & Not started \\
\bottomrule
\end{tabular}
\caption{Game mechanics documents.}
\label{tab:game-systems}
\end{table}

\subsection{Presentation}

\begin{table}[h]
\centering
\begin{tabular}{@{}llp{5.5cm}l@{}}
\toprule
\textbf{Doc} & \textbf{System} & \textbf{Scope} & \textbf{Status} \\
\midrule
13 & Graphics, Lighting \& Art Pipeline & 2D raycasting, normal maps, shadow stealth, art phases & Designed \\
\bottomrule
\end{tabular}
\caption{Presentation layer documents.}
\label{tab:presentation}
\end{table}

\subsection{Content Design}

Game content that populates the simulation---enemies, missions, levels,
and narrative.  These documents are relatively independent of the
simulation architecture.

\begin{table}[h]
\centering
\begin{tabular}{@{}llp{5.5cm}l@{}}
\toprule
\textbf{Doc} & \textbf{System} & \textbf{Scope} & \textbf{Status} \\
\midrule
14 & Creatures \& AI Behaviors & Swarms, predators, genetic soldiers, AI design & Brainstorm \\
15 & Missions \& Objectives & Mission types, objectives, extraction/sabotage & Brainstorm \\
16 & Ship Layouts \& Level Design & Map design, room types, level editor & Brainstorm \\
17 & Narrative \& Lore & Satirical 2087 world, Watergate/Princes plot, factions & Designed \\
\bottomrule
\end{tabular}
\caption{Content design documents.}
\label{tab:content}
\end{table}

%----------------------------------------------------------------------
\section{Language Pipeline}
\label{sec:language}
%----------------------------------------------------------------------

\subsection{Prototyping Phase: Python}

The simulation is prototyped in \textbf{Python with NumPy}.  This
enables rapid iteration on game rules, physics parameters, and system
interactions before committing to a compiled language.  All propagation
systems (explosions, decompression, smoke) have been designed and
validated as NumPy array operations.

\subsection{Production Phase: C++ with Python Bindings}

Once core systems are validated, the simulation moves to \textbf{C++}
for performance (pathfinding, propagation, and bulk array operations at
thousands of steps per second).  The ML training pipeline remains in
\textbf{Python} (PyTorch).  The two are connected via \textbf{pybind11}
(or nanobind), exposing the C++ simulation as an importable Python
module with zero-copy array sharing.

\subsection{Architecture}

\begin{verbatim}
  +------------------+       +-------------------+
  |  C++ Simulation  | <---> |  pybind11 bindings|
  |  (core logic,    |       |  (thin wrapper,   |
  |   array data)    |       |   zero-copy NumPy)|
  +------------------+       +-------------------+
         |                            |
         |  same source               |  import sim
         v                            v
  +------------------+       +-------------------+
  |  Engine Plugin   |       |  Python ML Script |
  |  (game build)    |       |  (training loop)  |
  +------------------+       +-------------------+
\end{verbatim}

\subsection{Zero-Copy Array Sharing}

The C++ simulation class owns the raw array data.  pybind11 wraps
these buffers as NumPy arrays \emph{without copying}---both languages
point at the same memory.  From Python, accessing
\texttt{env.state.pressure} returns a standard NumPy array that can be
fed directly into a neural network.

A typical training loop:

\begin{verbatim}
  import sim          # compiled C++ module
  env = sim.BoardingEnv(map="frigate_01")

  obs = env.reset()   # returns dict of NumPy arrays (zero-copy)
  for step in range(max_steps):
      action = agent.act(obs)
      obs, reward, done, info = env.step(action)
      agent.learn(obs, reward, done)
\end{verbatim}

This follows the standard Gymnasium interface, making it compatible with
existing RL libraries (Stable Baselines3, RLlib, CleanRL, etc.).

\subsection{Build Targets}

The C++ simulation compiles into two targets from the \textbf{same
source code}:

\begin{description}[style=nextline]
    \item[Engine Plugin]
        Compiled as a module within the game engine project.  The
        presentation layer calls directly into the simulation via C++
        APIs.
    \item[Python Module (\texttt{.so} / \texttt{.pyd})]
        Built with CMake + pybind11.  Produces an importable Python
        package for headless training, evaluation, and analysis.
\end{description}

\noindent This ensures that the game and the training environment run
\emph{identical} simulation logic---no translation bugs, no drift
between the two.

\subsection{Precedent}

This pattern is standard in high-performance RL research.  Notable
examples include DeepMind's MuJoCo bindings, OpenAI Gym's Atari
environments, and PettingZoo's multi-agent environments---all C++
simulations with thin Python wrappers.

%----------------------------------------------------------------------
\section{Neural Network Considerations}
\label{sec:nn}
%----------------------------------------------------------------------

The array-based state representation is designed to be directly
consumable by convolutional neural networks, following the stacked
feature-plane approach used by AlphaStar (StarCraft~II).

\begin{itemize}
    \item Each cached array becomes one (or more) input channels.
    \item For training, the fine grid can be downsampled to coarse
          resolution if full resolution proves unnecessary.
    \item The headless simulation enables massively parallel self-play
          without rendering overhead.
    \item Action space definition and reward shaping are deferred to a
          future document.
\end{itemize}

\medskip
\noindent\textbf{Design principle:}  No architectural decision should
foreclose the option of training a neural network agent.  The grid
representation, clean state serialization, and headless simulation
capability are all in service of this goal.

%----------------------------------------------------------------------
\section{Open Questions \& Next Steps}
%----------------------------------------------------------------------

\begin{enumerate}
    \item \textbf{Engine choice:} Review engine candidates with Idir
          and select one (Section~\ref{sec:engine}).
    \item \textbf{Presentation layer integration:} Once the engine is
          chosen---how does the presentation layer poll or subscribe to
          simulation state?  Tile rendering approach?  Sprite/animation
          pipeline?
    \item \textbf{Convert brainstorm notes to design documents:}
          Prioritize Documents~01--07 (environmental simulation), then
          08--12 (game mechanics).
    \item \textbf{Prototype:} Build the Python/NumPy simulation
          prototype with core environmental systems.
    \item \textbf{Establish shared repository} and document conventions.
    \item \textbf{Network replication strategy} if multiplayer is
          pursued.
\end{enumerate}

\end{document}
